\section{Related Work}

\textbf{Process Reward Model}
Process reward models (PRMs) provide fine-grained reward signals at each
reasoning step and have substantially improved LLM performance on complex
multi-step tasks such as mathematical problem solving and code generation
\citep{uesato2022solving, lightman2023lets}.
Compared with outcome reward models that supervise only the final answer, PRMs
can localize intermediate errors more precisely, improve temporal credit
assignment, and mitigate the alignment gap where the answer is correct but the
reasoning process is flawed.
As a result, PRMs have become a core technique for test-time scaling, such as
Best-of-\(N\) \citep{lin2025toolprmfinegrainedinferencescaling}, and
reinforcement learning
\citep{DBLP:conf/iclr/LiL25a, xi2025agentprmprocessrewardmodels, mit2025process}.
Their main bottleneck is that high-quality process annotations are difficult to
obtain: manual labeling is expensive, limited in scale, and often subjective.

\textbf{Automated Annotation and Credit Redistribution}
Representative automated process-labeling methods include Math-Shepherd
\citep{wang2024math} and related Monte Carlo (MC)-based step annotation
approaches
\citep{zhang2025process, zhang-etal-2025-lessons, chen-etal-2025-better-process, ding2025scan}.
As LLM capabilities have advanced, LLM-as-a-judge methods have also become
popular \citep{zeng2025versaprm, duan2025efficient}.
MC labels are often noisy and require large rollout budgets, while judge-based
labels can inherit model biases and remain sensitive to prompts and scoring
criteria.
Recent work therefore combines MC labels with LLM judges to balance quality and
scalability \citep{wang2025athena, zhang-etal-2025-lessons}, but extra rollouts
and judge calls remain costly.
Credit-assignment methods for delayed rewards provide a complementary view:
RUDDER redistributes terminal returns through contribution analysis while
preserving the trajectory objective in a return-equivalent decision process
\citep{arjona2019rudder}.
Shapley-based model explanations offer a principled allocation rule, but their
outputs depend on how the value function, background distribution, and
explanation context are operationalized \citep{pmlr-v119-sundararajan20b}.
This motivates an internal differentiable scorer and Aumann--Shapley-style
attribution rather than external distributional Shapley estimates.

In contrast, our method treats reasoning trajectories as ordered structures with
explicit causal constraints and derives step-level credit assignment directly
from sparse outcome signals.
Rather than using ordered discrete value differences alone, we use this
cooperative-game view to motivate a differentiable trajectory scorer and an
Aumann--Shapley-style local credit signal.
This mechanism requires neither intermediate annotations nor external judges.
The local branch performs trajectory-specific credit assignment, while the
global branch provides a cross-trajectory reward-reshaping prior through
nearest-neighbor statistics.
In doing so, it turns sparse outcome supervision into reusable step-level
supervision while reducing dependence on expensive external annotation
pipelines and mitigating judge-induced bias and label instability.
